\documentclass[11pt]{article}

\usepackage[margin=1in]{geometry}
\usepackage{amsmath,amssymb,amsthm,mathtools}
\usepackage[hidelinks]{hyperref}
\usepackage[nameinlink,capitalize]{cleveref}

\newtheorem{theorem}{Theorem}[section]
\newtheorem{lemma}[theorem]{Lemma}
\newtheorem{proposition}[theorem]{Proposition}
\newtheorem{corollary}[theorem]{Corollary}
\theoremstyle{remark}
\newtheorem{remark}[theorem]{Remark}

\newcommand{\cC}{\mathcal C}
\newcommand{\cU}{\mathcal U}
\newcommand{\cV}{\mathcal V}

\title{Counting convex subsets in iterated order-type blow-ups}
\author{Authors omitted for circulation}
\date{Draft of August 13, 2026}

\begin{document}
\maketitle

\begin{abstract}
For a planar point set $P$ in general position, let $v(P)$ be the number
of its subsets in convex position, and let
\[
 f(N)=\min\{v(P):|P|=N\}.
\]
Erd\H{o}s and Hammer asked for the asymptotic behaviour of $f(N)$; a
standard counting estimate applied to the classical construction gives
$\log_2 f(N)\le (1+o(1))(\log_2N)^2$.
We prove
\[
 \limsup_{N\to\infty}
 \frac{\log_2 f(N)}{(\log_2N)^2}\le\frac12.
\]
The construction is a directionally specified iterated order-type blow-up.
We give exact substitution formulas for the numbers of caps, cups, and
convex subsets.  If a fixed $r$-point template has largest cap and cup of
sizes $a$ and $b$, respectively, its iterates have normalized coefficient
$(a+b-2)/(2\log_2r)$.  Balanced instances of the classical cup--cap
construction make this coefficient tend to $1/2$.  The cup--cap theorem
shows that $1/2$ is optimal among fixed-template iterations.  More strongly,
we prove that every point set obtained by arbitrary recursive strong
separations has logarithmic convex-subset count at least
$(\log_2N)^2/2-O((\log N)^{3/2})$.  Thus the
optimal universal coefficient $1/2$ is attained within this nonstationary
class.
\end{abstract}

\section{Introduction}

For a finite set $P\subset\mathbb R^2$ in general position, call a subset
of $P$ \emph{convex} if all of its points are vertices of its convex hull.
Subsets of size at most two are included; changing this convention, or
including the empty set, has no effect on any normalized asymptotic below.
Write $v(P)$ for the number of convex subsets of $P$.
Throughout the paper, all logarithms are base two.

Erd\H{o}s~\cite{Erdos1978} recorded a question posed with Hammer: determine
or estimate
\[
 f(N)=\min_{|P|=N}v(P),
\]
where the minimum is over planar point sets in general position.  The
classical Erd\H{o}s--Szekeres construction and the convex-polygon theorem
show that $f(N)=2^{\Theta((\log N)^2)}$; see, for example, the survey of
Morris and Soltan~\cite{MorrisSoltan2000}.  The open problem asks in
particular whether $\log f(N)/(\log N)^2$ has a limit.

For completeness, the standard coefficient-$1$ upper estimate is immediate.
Take $k=\lceil\log N\rceil+2$ and retain $N$ points from the classical
$2^{k-2}$-point construction with no convex $k$-set.  Then
\[
 v(P)\le\sum_{j<k}\binom Nj\le kN^{k-1}
      =2^{(1+o(1))(\log N)^2}.
\]

For context, Suk's bound $ES(k)=2^{k+o(k)}$~\cite{Suk2017} gives the
following standard lower estimate.  Put $t=\lfloor N^{1/2}\rfloor$ and
$k=(1-o(1))\log t$, so every $t$-point set contains a convex $k$-set.
Double-counting pairs $(K,T)$ with $K\subseteq T\subseteq P$, $|K|=k$,
$|T|=t$, gives
\[
 \#\{\text{convex $k$-subsets of $P$}\}
 \ge \frac{\binom Nt}{\binom{N-k}{t-k}}
 =\frac{\binom Nk}{\binom tk}.
\]
Consequently
\[
 \liminf_{N\to\infty}
 \frac{\log f(N)}{(\log N)^2}\ge\frac14.
\]
Together with \cref{thm:main}, this gives the rigorous window
\[
 \frac14\le\liminf\frac{\log f(N)}{(\log N)^2}
 \le\limsup\frac{\log f(N)}{(\log N)^2}\le\frac12.
\]

Our result improves the coefficient in the classical upper bound.

\begin{theorem}\label{thm:main}
One has
\[
 \limsup_{N\to\infty}
 \frac{\log f(N)}{(\log N)^2}\le\frac12.
\]
\end{theorem}

The construction is a special realization of an order-type blow-up.
General blow-ups and their iteration appear in
Han--Kohayakawa--Sales--Stagni~\cite{HanEtAl2019}; almost-vertical blow-ups
tailored to the Erd\H{o}s--Szekeres problem appear in
Baek--Balko~\cite{BaekBalko2026}.  Exact weighted identities for numbers of
convex polygons are also known~\cite{HuemerEtAl2022}.  The feature used here
is a prescribed orientation for triples meeting exactly two clusters.
Those orientations give exact substitution identities for the three
statistics in \cref{lem:composition}; iterating them is what produces the
geometric coefficient $1/2$.  The same base-normalized coefficient occurs
in Sz\'ekely's graph-theoretic analogue~\cite{Szekely1984}, but that argument
does not directly yield the geometric statement here.

We also determine the asymptotic answer inside the much larger class of
all recursively strongly separated point sets; see \cref{thm:strong-class}.
The lower bound there uses a max-endpoint recurrence and a multiscale reset
argument.  It does not extend to arbitrary order types, so the existence
and value of the limit in the Erd\H{o}s--Hammer problem remain open.

\section{A directional vertical composition}

After a rotation, order every point set by increasing $x$-coordinate.  For
three points $p<q<r$, let $\chi(p,q,r)\in\{-,+\}$ denote their orientation.
A \emph{cap} is a nonempty set all of whose increasing triples have sign
$-$, and a \emph{cup} is defined with sign $+$.  Singletons and pairs are
both caps and cups.

A sufficiently strong shear $(x,y)\mapsto(x,y+Mx)$ preserves orientations
and makes the $y$-coordinates increase in the same order as the
$x$-coordinates.  We shall therefore assume both coordinates increase.

Let $S=(s_1,\ldots,s_r)$ and $Q=(q_1,\ldots,q_n)$ have this property, with
$s_i=(X_i,Y_i)$ and $q_j=(x_j,y_j)$.  For $\varepsilon>0$, replace $s_i$ by
the block
\[
 Q_i=\bigl\{(X_i+\varepsilon^2x_j,
              Y_i+\varepsilon y_j):1\le j\le n\bigr\}.       \tag{2.1}
\]
For all sufficiently small $\varepsilon$, denote the resulting point set
by $S[Q]$.

\begin{lemma}[orientation rules]\label{lem:orientation}
For sufficiently small positive $\varepsilon$, the following hold in
$S[Q]$.
\begin{enumerate}
\item A triple in one block has its orientation from $Q$.
\item A triple in three different blocks has its orientation from $S$.
\item If the first two points of an increasing triple are in one block,
      its orientation is $-$.
\item If the last two points are in one block, its orientation is $+$.
\end{enumerate}
The parameter $\varepsilon$ may be chosen rational, and $S[Q]$ again has
both coordinates strictly increasing.  If $S$ and $Q$ have rational
coordinates, then so does $S[Q]$.
\end{lemma}

\begin{proof}
The first assertion follows by scaling within a block, and the second by
continuity from the corresponding triple of cluster centres.  If $p<q$
are in one block and $z$ is in a later block, then
\[
 q-p=(\varepsilon^2\Delta x,\varepsilon\Delta y),
 \qquad \Delta x,\Delta y>0.
\]
In the determinant $\det(q-p,z-p)$, the term
$-\varepsilon\Delta y\,(z_x-p_x)$ dominates, proving the third assertion.
The fourth is the reflected calculation.

Only finitely many strict coordinate inequalities and nonzero determinant
conditions are involved.  Hence they hold throughout some interval
$0<\varepsilon<\varepsilon_0$; choose a rational value there.  Shrinking
$\varepsilon_0$ if necessary also preserves the coordinate order.
In an iteration, this choice is made afresh at every finite stage; no
uniform perturbation parameter is asserted or needed.
\end{proof}

For a point set $R$, let $c_j(R),u_j(R),v_j(R)$ denote the numbers of
$j$-point caps, cups, and convex subsets.  Put
\[
 C(R)=\sum_{j\ge1}c_j(R),\qquad
 U(R)=\sum_{j\ge1}u_j(R),\qquad
 W(R)=\sum_{j\ge1}v_j(R).                         \tag{2.2}
\]

\begin{lemma}[exact composition count]\label{lem:composition}
If $|S|=r$ and $|Q|=n$, then
\begin{align}
 C(S[Q])&=C(Q)\sum_{j\ge1}c_j(S)n^{j-1},\tag{2.3}\label{eq:Ccomp}\\
 U(S[Q])&=U(Q)\sum_{j\ge1}u_j(S)n^{j-1},\tag{2.4}\label{eq:Ucomp}\\
 W(S[Q])&=rW(Q)+C(Q)U(Q)
                 \sum_{j\ge2}v_j(S)n^{j-2}.\tag{2.5}\label{eq:Wcomp}
\end{align}
\end{lemma}

\begin{proof}
Consider a cap meeting at least two blocks.  Only its first occupied block
may contain more than one point: two points in any later block, together
with a point in an earlier block, form a positive triple by
\cref{lem:orientation}.  The first block contributes an arbitrary nonempty
cap of $Q$, every other occupied block contributes one arbitrary point, and
the occupied block indices form a cap of $S$.  The orientation rules also
show these conditions are sufficient.  Summing over macro-caps proves
\eqref{eq:Ccomp}; reflection proves \eqref{eq:Ucomp}.

Now let $X\subseteq S[Q]$ be convex and meet at least two blocks.  Its upper
and lower hull chains are respectively a cap and a cup with common
endpoints.  The lower hull contains at most one point from the first
occupied block: two such points followed by any selected point in a later
block would have sign $+$ along the lower chain, contradicting
\cref{lem:orientation}.  Since every selected point is a hull vertex, all
remaining points of the first block lie on the upper chain.  Thus the first
occupied block intersects $X$ in a cap.  The reflected argument makes the
last-block intersection a cup.  We claim every intermediate block contains
at most one point.

Indeed, suppose selected points $a<b_1<b_2<c$ lie in that order, with
$b_1,b_2$ in one intermediate block and $a,c$ in earlier and later blocks.
For small $\varepsilon$, the points $b_1,b_2$ lie on the same side of the
line $ac$: this is true of their common cluster centre, and all relevant
determinants are strict.  If they lie above $ac$, convexity would put both
on the upper chain from $a$ to $c$, forcing
$\chi(a,b_1,b_2)=-$; if they lie below, both lie on the lower chain,
forcing $\chi(b_1,b_2,c)=+$.  Both conclusions contradict
\cref{lem:orientation}.

Choose one representative from each occupied endpoint block and the unique
point in every occupied intermediate block.  This is a subset of the
convex set $X$, hence is convex, and its order type is the corresponding
subset of $S$.  Thus the occupied block indices form a convex subset of
$S$.

Conversely, fix a convex $j$-subset of $S$, $j\ge2$.  Choose a nonempty cap
of $Q$ in its first block, a nonempty cup of $Q$ in its last block, and one
point in each intermediate block.  Along the upper macro-hull, insert the
entire first-block cap and the rightmost point of the last-block cup; the
orientation rules make the result a strict cap.  The reflected construction
along the lower macro-hull is a strict cup containing the remaining chosen
points.  These two chains share their endpoints, lie on opposite sides of
their endpoint chord, and do not cross.  Their union is therefore the
boundary of a convex polygon.  The choices are unique and contribute
$C(Q)U(Q)n^{j-2}$.  Sets contained in a single block contribute $rW(Q)$,
which proves \eqref{eq:Wcomp}.
\end{proof}

\section{Iteration of a fixed template}

Fix a template $S$ of size $r\ge2$.  Let its largest cap and cup have sizes
$a$ and $b$, and define
\[
 Q_0=\{\text{one point}\},\qquad Q_d=S[Q_{d-1}].              \tag{3.1}
\]
Thus $|Q_d|=r^d$.

\begin{proposition}\label{prop:template-rate}
The iterates satisfy
\[
 \lim_{d\to\infty}
 \frac{\log W(Q_d)}{(\log|Q_d|)^2}
 =\frac{a+b-2}{2\log r}.                          \tag{3.2}
\]
\end{proposition}

\begin{proof}
Define the fixed polynomials
\[
 A(z)=\sum_{j\ge1}c_j(S)z^{j-1},\quad
 B(z)=\sum_{j\ge1}u_j(S)z^{j-1},\quad
 D(z)=\sum_{j\ge2}v_j(S)z^{j-2}.                 \tag{3.3}
\]
Their degrees are $a-1$, $b-1$, and at most $r-2$, respectively; each has
positive leading coefficient, and $D$ is nonzero because every pair is
convex.  From \cref{lem:composition},
\[
 C(Q_d)=C(Q_{d-1})A(r^{d-1}),\qquad
 U(Q_d)=U(Q_{d-1})B(r^{d-1}).                     \tag{3.4}\label{eq:CU-recurrence}
\]
Since a fixed positive polynomial $P$ of degree $e$ satisfies
$\log P(r^t)=et\log r+O_S(1)$ uniformly for $t\ge0$, summing
\eqref{eq:CU-recurrence} over the levels gives
\begin{align}
 \log C(Q_d)&=(a-1)\binom d2\log r+O_S(d),\tag{3.5}\label{eq:Citer}\\
 \log U(Q_d)&=(b-1)\binom d2\log r+O_S(d).\tag{3.6}\label{eq:Uiter}
\end{align}

The third recurrence unrolls exactly as
\[
 W(Q_d)=r^dW(Q_0)+
 \sum_{s=1}^d r^{d-s}C(Q_{s-1})U(Q_{s-1})D(r^{s-1}).          \tag{3.7}
\]
By \eqref{eq:Citer}--\eqref{eq:Uiter}, the logarithm of the $s$th summand is
\[
 (a+b-2)\binom{s-1}{2}\log r+O_S(s)+O_S(d-s).                 \tag{3.8}
\]
There are only $d$ summands, and the last one has the displayed quadratic
term with $s=d$.  More explicitly, since
$D(r^{d-1})\ge v_2(S)>0$, the last summand gives
\[
 W(Q_d)\ge v_2(S)C(Q_{d-1})U(Q_{d-1}).
\]
Hence upper-bounding the sum by $d$ times its largest quadratic-scale bound
and using this explicit lower bound yields
\[
 \log W(Q_d)=\frac{a+b-2}{2}d^2\log r+O_S(d).                 \tag{3.9}
\]
Division by $(d\log r)^2$ proves the proposition.
\end{proof}

\section{Balanced cup--cap templates}

We first record the geometric operation used in the classical recursive
cup--cap construction.  Write $A\prec B$ when internal orientations are
preserved and every increasing triple meeting both sides has sign $-$ if
its first two points lie in $A$, and sign $+$ if its last two points lie in
$B$.

\begin{lemma}[strong separation]\label{lem:strong-glue}
Given rational point sets $A$ and $B$ in general position, they have a
rational realization $A\prec B$ in general position.  Moreover, a cap in
$A\prec B$ meeting both sides consists of a cap of $A$ and one point of
$B$, while a cup meeting both sides consists of one point of $A$ and a cup
of $B$.
\end{lemma}

\begin{proof}
First shear both sets so that their two coordinates increase.  Place a
copy of $A$ in the block
$(\varepsilon^2x,\varepsilon y)$ near $(0,0)$ and a copy of $B$ in the
block $(1+\varepsilon^2x,1+\varepsilon y)$ near $(1,1)$.  The determinant
calculation from \cref{lem:orientation} gives the two required mixed signs
for all sufficiently small $\varepsilon>0$.  A rational $\varepsilon$ may
be selected below all finitely many strict thresholds.  The cap and cup
classifications follow immediately: two points of the later block forbid a
cap, while two points of the earlier block forbid a cup; the mixed signs
also prove the converse.
\end{proof}

Now let $T_{m,0}$ and $T_{m,m}$ be singletons, and for $0<i<m$ form
\[
 T_{m,i}=T_{m-1,i-1}\prec T_{m-1,i}.             \tag{4.1}\label{eq:pascal-glue}
\]
using \cref{lem:strong-glue}.  These are the standard cup--cap examples.

\begin{lemma}\label{lem:pascal}
The cell $T_{m,i}$ has $\binom mi$ points, largest cap of size $i+1$, and
largest cup of size $m-i+1$.
\end{lemma}

\begin{proof}
The size recurrence is Pascal's identity.  For caps, the rule in
\eqref{eq:pascal-glue} gives the larger of the left-child cap size plus one and the
right-child cap size.  Induction gives $i+1$ in either case.  Reflection
gives the cup statement.  The boundary convention is consistent because
the boundary cells are singletons.
\end{proof}

For $k\ge3$, take the central cell
\[
 S_k=T_{2k-4,k-2}.                                \tag{4.2}
\]
Then
\[
 r_k=|S_k|=\binom{2k-4}{k-2},\qquad a_k=b_k=k-1. \tag{4.3}
\]

\begin{proof}[Proof of \cref{thm:main}]
Fix $k$ and iterate $S_k$.  By \cref{prop:template-rate}, configurations
of size $r_k^d$ satisfy
\[
 \frac{\log W(Q_d)}{(\log|Q_d|)^2}
 \longrightarrow
 \frac{k-2}{\log\binom{2k-4}{k-2}}.              \tag{4.4}
\]
For arbitrary $N$, take the least $d$ with $r_k^d\ge N$ and retain any
$N$ points.  Deletion creates no new convex subsets, while
$r_k^d<r_kN$, so the normalized coefficient is unchanged.  Consequently
\[
 \limsup_{N\to\infty}
 \frac{\log f(N)}{(\log N)^2}
 \le \frac{k-2}{\log\binom{2k-4}{k-2}}.          \tag{4.5}
\]
Stirling's formula gives
$\log\binom{2k-4}{k-2}=2k-O(\log k)$.  Letting $k\to\infty$ proves the
theorem.
\end{proof}

The same calculation, together with the cup--cap theorem of Erd\H{o}s and
Szekeres~\cite{ErdosSzekeres1935}, identifies the limitation of
fixed-template iteration.

\begin{proposition}\label{prop:barrier}
For every fixed template $S$, the coefficient in
\cref{prop:template-rate} is at least $1/2$.
\end{proposition}

\begin{proof}
If the largest cap and cup have sizes $a$ and $b$, the cup--cap theorem
applied with forbidden sizes $a+1$ and $b+1$ gives
\[
 r\le\binom{a+b-2}{a-1}\le2^{a+b-2}.
\]
Thus $(a+b-2)/(2\log r)\ge1/2$.
\end{proof}

\section{Optimality for arbitrary strong decompositions}

Call a point set \emph{strongly decomposable} if it is a singleton, or if
it has a realization $A\prec B$ in which $A$ and $B$ are strongly
decomposable.  Equivalently, it is represented by an ordered full binary
tree whose internal nodes are the operation in \cref{lem:strong-glue}.
Let
\[
 g(N)=\min\{W(P): |P|=N,\ P\text{ strongly decomposable}\}.
\]

\begin{theorem}\label{thm:strong-class}
Every $N$-point strongly decomposable set $P$ satisfies
\begin{equation}
 \log W(P)\ge \frac12(\log N)^2-O((\log N)^{3/2}),
 \label{eq:strong-lower}
\end{equation}
with an absolute implicit constant.  Consequently
\begin{equation}
 \lim_{N\to\infty}\frac{\log g(N)}{(\log N)^2}=\frac12.
 \label{eq:strong-limit}
\end{equation}
\end{theorem}

We need two preparatory estimates.  At a node $T=A\prec B$, write
$a=|A|$ and $b=|B|$.  The strong-separation classification gives
\begin{align}
 C(T)&=C(B)+(b+1)C(A),\label{eq:tree-C}\\
 U(T)&=U(A)+(a+1)U(B).\label{eq:tree-U}
\end{align}

\begin{lemma}[cap--cup product]\label{lem:product-mass}
Every $s$-point strongly decomposable set $S$ satisfies
\[
 \log C(S)+\log U(S)\ge\frac12(\log s)^2-\log s.
\]
\end{lemma}

\begin{proof}
Put $R(S)=\sqrt{C(S)U(S)}$.  The Cauchy--Schwarz inequality
$\sqrt{(p+r)(q+t)}\ge\sqrt{pq}+\sqrt{rt}$ and
\eqref{eq:tree-C}--\eqref{eq:tree-U} give
\[
 R(A\prec B)\ge \sqrt{b+1}\,R(A)+\sqrt{a+1}\,R(B).
\]
Starting at $S$, follow a larger child until reaching a leaf.  If the
current subtree has $m_i$ points, the followed child has
$m_{i+1}=m_i-s_i$ points, where $1\le s_i\le m_i/2$.  Iterating the last
display gives
\begin{equation}
 \log R(S)\ge\frac12\sum_i\log(s_i+1).
 \label{eq:R-path}
\end{equation}
Let $t_i=\log m_i$ and $d_i=t_i-t_{i+1}$.  Since $0<d_i\le1$, convexity of
$x^d$ on $0\le d\le1$, with $x=m_i/2$, shows
\[
 (m_i/2)^{d_i}\le1+d_i(m_i/2-1)\le s_i+1;
\]
here we used $d_i=-\log(1-s_i/m_i)\le2s_i/m_i$.  Consequently
\[
 \sum_i\log(s_i+1)
 \ge\sum_i d_i(t_i-1)
 \ge\frac12(\log s)^2-\log s,
\]
because $\sum_i d_i=\log s$ and
$\sum_i d_it_i=((\log s)^2+\sum_i d_i^2)/2$.  Combining this with
\eqref{eq:R-path} gives
$\log R(S)\ge\frac14(\log s)^2-\frac12\log s$; doubling proves
the claim.
\end{proof}

We next refine the count by endpoints.  For $p<q$, let $c(p,q)$ and
$u(p,q)$ be the numbers of caps and cups with endpoints $p,q$.  Let $x(p)$
count caps with left endpoint $p$, including the singleton, and let $y(q)$
count cups with right endpoint $q$, again including the singleton.  Set
\[
 X=\max_p x(p),\qquad Y=\max_q y(q),\qquad
 M=\max\bigl(1,\max_{p<q}c(p,q)u(p,q)\bigr).
\]
For a leaf, $(X,Y,M)=(1,1,1)$.  Directly from strong separation,
\begin{align}
 X(A\prec B)&=\max\{(b+1)X(A),X(B)\},\label{eq:X-rec}\\
 Y(A\prec B)&=\max\{Y(A),(a+1)Y(B)\},\label{eq:Y-rec}\\
 M(A\prec B)&=\max\{M(A),M(B),X(A)Y(B)\}.\label{eq:M-rec}
\end{align}
Indeed, for crossing endpoints $p\in A$, $q\in B$, one has exactly
$c(p,q)=x_A(p)$ and $u(p,q)=y_B(q)$.

Every nonsingleton convex set has a unique upper cap and lower cup with the
same endpoints, and conversely every such pair forms a convex set.  Hence
\begin{equation}
 W(S)=|S|+\sum_{p<q}c(p,q)u(p,q),
 \qquad M(S)\le W(S)\le2|S|^2M(S).
 \label{eq:WM}
\end{equation}
Also $x(p)=1+\sum_{q>p}c(p,q)$, and every two-point cup is available, so
\begin{equation}
 X(S)\le |S|M(S),\qquad Y(S)\le |S|M(S).
 \label{eq:XY-ceiling}
\end{equation}
Finally, $C(S)\le |S|X(S)$ and $U(S)\le |S|Y(S)$.  Thus
\cref{lem:product-mass} implies, writing lower-case letters for base-two
logarithms,
\begin{equation}
 x(S)+y(S)\ge\frac12(\log|S|)^2-3\log|S|.
 \label{eq:radial}
\end{equation}

\begin{proof}[Proof of \cref{thm:strong-class}]
It is enough to consider all sufficiently large $N$; increase the absolute
constant to absorb the remaining values.  Put $L=\log N$,
$R=\lceil\sqrt L\rceil$, and $\lambda=\log L$.  In the decomposition tree
of $P$, start at the root and repeatedly follow a larger child.  If $n_i$
is the current size and $s_i$ the discarded sibling size, stop when first
\[
 n_i<N/2^{4R}.
\]
Call a level large when $s_i\ge n_i/L^2$.

Suppose first that fewer than $R$ levels are large.  A nonlarge level loses
at most $2/L^2$ bits of size, since
$-\log(1-s_i/n_i)\le2/L^2$, while every level loses at most one bit.
The total loss before stopping exceeds $4R$, so there are more than
$\frac32RL^2$ nonlarge levels.  At least half have their discarded sibling
on the same side of the path.  Fix one leaf in each such sibling and one
leaf below the stopping node.  Choosing an arbitrary subset of the fixed
sibling leaves, together with the terminal leaf, gives distinct caps if
the siblings lie on the right, and distinct cups if they lie on the left.
This follows inductively upward from the mixed signs in
\cref{lem:strong-glue}.  Therefore
\[
 \log W(P)>\frac34RL^2,
\]
which is stronger than \eqref{eq:strong-lower}.

Now suppose there are at least $R$ large levels.  At each of $R$ selected
large nodes, both children have at least $2^{L-\Delta}$ leaves, where
\[
 \Delta=4R+2\lambda+1.
\]
Indeed, the followed child has at least $n_i/2$, the sibling at least
$n_i/L^2$, and $n_i\ge N/2^{4R}$ before the stop.  Since
$L-\Delta>3$ for all sufficiently large $L$ and
$u\mapsto u^2/2-3u$ is then increasing in the relevant range,
\eqref{eq:radial} gives uniformly for each such child
\[
 x+y\ge F:=\frac12(L-\Delta)^2-3L.
\]

Let $\mu=\log M(P)$.  Equations \eqref{eq:X-rec}--\eqref{eq:M-rec} make
$X,Y,M$ nondecreasing from a child to its parent.  By
\eqref{eq:XY-ceiling}, every coordinate of every subtree is at most
$\mu+L$.  If $\mu\ge F-L$, the desired estimate follows immediately.
Otherwise put
\[
 D=F-\mu>L,\qquad \ell=D-L>0.
\]
Every coordinate of every child at a selected node is at least $\ell$, by
\eqref{eq:radial} and the ceiling $\mu+L$.

Order the $R$ selected nodes from deepest to highest.  At the deepest node
$A\prec B$, the crossing term in \eqref{eq:M-rec} gives
$x_A+y_B\le\mu$.  Together with the two radial bounds
\eqref{eq:radial}, this implies
\[
 x_B,y_A\ge2D-L.
\]
Thus both coordinates of the parent are at least $h_0:=2D-L$.

At a later selected node, let $H$ be the child containing the preceding
selected node and $Q$ the opposite child.  If $H$ is the left child, then
$x_H+y_Q\le\mu$ and $x_Q+y_Q\ge F$, whence
\[
 x_Q\ge x_H+D.
\]
The parent increases its $x$-coordinate by at least $D$ and preserves its
$y$-coordinate.  If $H$ is the right child, the reflected argument
increases $y$ by at least $D$ and preserves $x$.  Intermediate nodes cannot
erase either gain, by monotonicity.

Among the $R-1$ later attachments, one direction occurs $q_*$ times, where
$q_*\ge\lceil(R-1)/2\rceil$.  At its final occurrence the forward product
uses a path coordinate of at least $h_0+(q_*-1)D$ and an opposite-child
coordinate of at least $\ell$.  Therefore
\[
 \mu\ge(q_*+2)D-2L,
 \qquad
 \mu\ge\frac{q_*+2}{q_*+3}F-\frac{2L}{q_*+3}.
\]
Since $\Delta=O(\sqrt L)$, $q_*=\Omega(\sqrt L)$, and
$F=\frac12L^2-O(L^{3/2})$, both alternatives give
\[
 \mu\ge\frac12L^2-O(L^{3/2}).
\]
Now $W(P)\ge M(P)$, proving \eqref{eq:strong-lower}.

For the limit in \eqref{eq:strong-limit}, observe that every $Q_d$
constructed from the Pascal templates is strongly decomposable.  Indeed,
at every internal split of the macro strong tree, triples whose two
same-side points lie in one micro-block have the required sign by
\cref{lem:orientation}(3)--(4), while triples using distinct micro-blocks
inherit the sign of the macro split.  Induction therefore identifies the
iterate with the strong tree obtained by substituting the micro strong tree
at every macro leaf.  An induced subset is also strong after suppressing
empty and unary nodes.  The arbitrary-$N$ deletion argument in the proof of
\cref{thm:main} therefore gives
$\limsup\log g(N)/(\log N)^2\le1/2$, while the first part gives the
matching liminf.
\end{proof}

\section{Verification artifact}

The accompanying exact-arithmetic program independently checks the finite
combinatorics behind \cref{lem:composition}.  It constructs the abstract
composition signs, realizes them with rational coordinates, counts caps and
cups by a last-edge dynamic program, and counts convex subsets through
upper/lower endpoint chains.  A $9$-point composition is additionally
checked by enumerating all $2^9$ subsets.  On
$S=Q=T_{4,2}$, the substitution formulas and independent dynamic programs
both return
\[
 (C,U,W)=(14136,14136,441399)
\]
for the resulting $36$-point composition.  The artifact is described in
the supplementary README and is not used as a premise of the proof.  A
second exact-arithmetic program checks the one-reset inequalities used in
\cref{thm:strong-class} over finite boxes of logarithmic states.  It is an
algebraic smoke test, not a substitute for the multiscale argument.

\small
\begin{thebibliography}{HOPTV22}

\bibitem[BB26]{BaekBalko2026}
J.~Baek and M.~Balko.
The Erd\H{o}s--Szekeres conjecture revisited.
\emph{J. Combin. Theory Ser. A}, 222:106195, 2026.

\bibitem[Erd78]{Erdos1978}
P.~Erd\H{o}s.
Some more problems on elementary geometry.
\emph{Austral. Math. Soc. Gaz.}, 5(2):52--54, 1978.

\bibitem[ES35]{ErdosSzekeres1935}
P.~Erd\H{o}s and G.~Szekeres.
A combinatorial problem in geometry.
\emph{Compos. Math.}, 2:463--470, 1935.

\bibitem[HKSS19]{HanEtAl2019}
J.~Han, Y.~Kohayakawa, M.~T. Sales, and H.~Stagni.
Extremal and probabilistic results for order types.
In \emph{Proc. 30th ACM--SIAM Symposium on Discrete Algorithms},
pages 426--435. SIAM, 2019.

\bibitem[HOPTV22]{HuemerEtAl2022}
C.~Huemer, D.~Oliveros, P.~P\'erez-Lantero, F.~Torra, and B.~Vogtenhuber.
On weighted sums of numbers of convex polygons in point sets.
\emph{Discrete Comput. Geom.}, 68:448--476, 2022.

\bibitem[MS00]{MorrisSoltan2000}
W.~Morris and V.~Soltan.
The Erd\H{o}s--Szekeres problem on points in convex position---a survey.
\emph{Bull. Amer. Math. Soc.}, 37(4):437--458, 2000.

\bibitem[Suk17]{Suk2017}
A.~Suk.
On the Erd\H{o}s--Szekeres convex polygon problem.
\emph{J. Amer. Math. Soc.}, 30(4):1047--1053, 2017.

\bibitem[Sz84]{Szekely1984}
L.~A. Sz\'ekely.
On the number of homogeneous subgraphs of a graph.
\emph{Combinatorica}, 4:363--372, 1984.

\end{thebibliography}
\normalsize

\end{document}
